\clearpage\begingroup\NoSourceHead\SourceLeaf{517b}
\renewcommand{\thefootnote}{\arabic{footnote})}\setcounter{footnote}{11}
\ArithmeticHead{60}{L}{BOLZA, GAUSS UND DIE VARIATIONSRECHNUNG.}
\begin{center}\textls[100]{2) Beziehungen zum Unabhängigkeitssatz.}\end{center}
\noindent\hspace*{\parindent}9. Noch tiefer in den eigentlichen Kern der modernen Variationsrechnung
führen uns die \textsc{Gauss}schen Untersuchungen vom Art. 21 und 22.

Ist ein Extremalenfeld für das Integral
\begin{equation*}\tag{8}
J=\int F(x,y,x',y')\,dt
\end{equation*}
gegeben und eine Transversale $\mathfrak C_0$ desselben, bezeichnen ferner $p(x,y),q(x,y)$
die Richtungskosinus der durch den Punkt $x,y$ gehenden Extremale des Feldes
im Punkt $x,y$ und $V(x,y)$ das \textls[100]{Feldintegral}\footnote{d. h. den Wert des Integrals $J$ genommen entlang der durch den Punkt $x,y$ gehenden Feldextremale
von ihrem Schnittpunkt mit der Kurve $\mathfrak C_0$ aus bis zum Punkt $x,y$.} bezogen auf die Ausgangskurve
$\mathfrak C_0$, so gelten für die partiellen Ableitungen der Funktion $V$ die »\textsc{Hamilton}schen Formeln«\footnote{Als Spezialfall in den allgemeineren \textsc{Hamilton}schen Formeln (35) enthalten, siehe \textsc{Hamilton},
loc. cit. 9), First supplement, Bd. 16, Art. 3; vgl. auch \textsc{Kneser}, \textit{Lehrbuch} (1900), §§ 15, 19.}
\begin{equation*}\tag{24}
\frac{\partial V}{\partial x}=F_{x'}(x,y,p,q),\quad\frac{\partial V}{\partial y}=F_{y'}(x,y,p,q),
\end{equation*}
woraus sich einerseits die Integrabilitätsbedingung
\begin{equation*}\tag{25}
\frac{\partial F_{x'}(x,y,p,q)}{\partial y}=\frac{\partial F_{y'}(x,y,p,q)}{\partial x},
\end{equation*}
andererseits durch Elimination von $p,q$ die »\textsc{Jacobi-Hamilton}sche partielle
Differentialgleichung«
\begin{equation*}\tag{26}
\Phi\left(x,y,\frac{\partial V}{\partial x},\frac{\partial V}{\partial y}\right)=0
\end{equation*}
für die Funktion $V$ ergibt. Von hier aus gelangt man dann unmittelbar zum
\clearpage\SourceLeaf{518}
\ArithmeticHead{61}{R}{BEZIEHUNGEN ZUM UNABHÄNGIGKEITSSATZ.}
\noindent Hinlänglichkeitsbeweis für ein Extremum des Integrals $J$, sei es mittels der
\textsc{Weierstrass}schen Konstruktion, sei es mittels des \textsc{Hilbert}schen invarianten
Integrals
\begin{equation*}\tag{27}
J^*=\int[F_{x'}(x,y,p,q)\,dx+F_{y'}(x,y,p,q)\,dy].
\end{equation*}

Die ersten Keime dieser Theorie finden sich bereits bei \textsc{Gauss} (Art. 22)
sowie in der schon oben genannten, fast gleichzeitig mit den \textit{Disquisitiones}
erschienenen Arbeit von \textsc{Hamilton}.
\par\noindent\hspace*{\parindent}a) Die \textsc{Gauss}schen Formeln.
\par\noindent\hspace*{\parindent}10. \textsc{Gauss} beweist in Art. 22 auf Grund seiner allgemeinen Resultate
über den Übergang von einem System krummliniger Koordinaten $p,q$ zu einem
zweiten System $p',q'$ für den speziellen Fall der geodätischen Parallelkoordinaten
$r,\varphi$ die Formeln
\begin{equation*}\tag{28}
\frac{\partial r}{\partial p}=\sqrt E\cos(\omega-\psi),\quad\frac{\partial r}{\partial q}=\sqrt G\cos\psi,
\end{equation*}
\begin{equation*}\tag{29}
E\left(\frac{\partial r}{\partial q}\right)^2-2F\frac{\partial r}{\partial p}\frac{\partial r}{\partial q}+G\left(\frac{\partial r}{\partial p}\right)^2=EG-F^2.
\end{equation*}

Dabei ist $r$ der senkrechte geodätische Abstand von einer festen Ausgangskurve
$\mathfrak C_0$, also mit dem Feldintegral für das geodätische Integral (4) identisch;
$\psi$ ist der Winkel, welchen die durch den Punkt $p,q$ gehende geodätische
Linie der zu $\mathfrak C_0$ orthogonalen Schar im Punkte $p,q$ mit der $q$-Richtung bildet,
während $\omega$ den Winkel zwischen der $p$-Richtung und der $q$-Richtung im Punkt
$p,q$ bedeutet. Daraus geht hervor, dass die \textsc{Gauss}sche Formel (29) mit der
\textsc{Jacobi-Hamilton}schen partiellen Differentialgleichung (26) für das Integral (4)
identisch ist, wie schon F. \textsc{Minding} erkannt hat\begingroup\renewcommand{\thefootnote}{13a)}\footnote{F. \textsc{Minding}, \textit{De formae, in quam geometra britannicus Hamilton integralia mechanices analyticae
redegit, origine genuina}, Dorpat 1864, wieder abgedruckt Mathematische Annalen, Bd. 55 (1902),
S. 119—135; vgl. insbesondere S. 122 und S. 131.}\endgroup, während die Formeln (28)
die \textsc{Hamilton}schen Formeln (24) in unentwickelter Form darstellen. Um die
vollständige Übereinstimmung zwischen beiden herzustellen, braucht man nur
mit Hilfe der von \textsc{Gauss} in Art. 17 gegebenen Formeln in (28) die trigonometrischen
Funktionen von $\omega$ und $\psi$ durch die drei Fundamentalgrössen $E,F,G$
und die Differentiale $dp,dq$ entlang der geodätischen Linie durch den Punkt
$p,q$ auszudrücken. Man erhält so in Übereinstimmung mit (24)
\begin{equation*}\tag{30}
\frac{\partial r}{\partial p}=\frac{E\,dp+F\,dq}{ds},\quad\frac{\partial r}{\partial q}=\frac{F\,dp+G\,dq}{ds}.
\end{equation*}
\clearpage\SourceLeaf{519}
\ArithmeticHead{62}{L}{BOLZA, GAUSS UND DIE VARIATIONSRECHNUNG.}
\noindent\hspace*{\parindent}Explizite findet sich diese Formel (30) übrigens bei \textsc{Gauss} nicht.

Bei der ganzen Untersuchung ist stillschweigend vorausgesetzt, dass man
sich auf ein solches Stück der Fläche beschränkt, dass durch jeden Punkt
dieses Stücks nur eine geodätische Linie senkrecht zur Ausgangskurve $\mathfrak C_0$ gezogen
werden kann, sodass also die betrachtete Schar von geodätischen Linien
ein Feld bildet.
\par\clearpage\endgroup
